\documentclass[11pt]{article}
\usepackage{series}
\usepackage{booktabs}
\usepackage{array}
\usepackage{bm}

\DeclareMathOperator{\Tr}{Tr}
\newcommand{\cE}{\mathcal E}
\newcommand{\cL}{\mathcal L}
\newcommand{\cU}{\mathcal U}

\title{\bf Electromagnetic Helicity from an Induced Line Connection\\
{\large Exact Chern--Simons Identities and a Conditional MTT Encoding}}
\author{Peter Nero}
\date{Version 3, July 2026}

\begin{document}
\maketitle

\begin{abstract}
A smooth rank-one orthogonal projector on a Hermitian bundle defines a line
subbundle and an induced unitary connection.  This elementary construction
provides a precise setting in which Berry curvature, Abelian Chern--Simons
functionals, magnetic helicity, and Hopf invariants can be compared.  We derive
the induced curvature, including its projector term, and formulate helicity only
after specifying a global trivialization or a relative-helicity protocol.  The
exact slice balance is the usual electric--magnetic pairing plus an explicit
boundary flux.  There is no additional projector remainder when the electric
and magnetic fields are those of the full induced connection: projector
variation is already part of those fields.  A quantitative comparison theorem
bounds the error made when the Berry curvature is omitted and only the ambient
Abelian field is retained.  We also give a correctly scaled Riesz-projector
derivative estimate and state the Hopf normalization with all conventions
visible.  Modal Triplet Theory (MTT) can use this construction as a conditional
encoding once a selected source emits the Hermitian bundle, connection,
rank-one projector, physical field identification, and boundary data.  Current
MTT results do not yet establish that source theorem for the physical
electromagnetic sector.  The paper therefore proves an exact geometric
dictionary and a controlled comparison result, not a universal electromagnetic
prediction.
\end{abstract}

\tableofcontents

\section*{Revision note for Version 3}

\begin{description}[leftmargin=2.8cm,style=nextline]
\item[Supersedes.]
Version 2, DOI \href{https://doi.org/10.5281/zenodo.18261452}
{10.5281/zenodo.18261452}.

\item[Reason.]
Version 2 correctly retained the Berry term in the induced curvature, but it
then counted projector variation a second time as a remainder in the exact
helicity balance.  It also contracted a spatial two-form to define the
electric field, mixed two Hopf normalizations, and treated a candidate
rank-one coherent line as the selected physical electromagnetic sector.

\item[Resolution.]
Version 3 distinguishes the full spacetime curvature from its spatial and
electric parts, proves the boundary-aware balance for the full induced field,
and uses projector control only for the error made by an ambient-field
surrogate.  It declares the global gauge and topology domain, fixes the Hopf
and Riesz constants, and makes physical MTT sourcing a six-row conditional
contract.

\item[Retained content.]
The projected-connection construction and the fact that its curvature contains
a Berry/Grassmann term remain valid.  The revision changes the balance-law
bookkeeping and physical claim status, not that underlying geometry.

\item[Open boundary.]
Current MTT results do not yet select the physical photon line, identify the
induced curvature with the physical electromagnetic field, or transport the
required dynamics.  Those source rows remain explicit conditions rather than
results of this paper.
\end{description}

\section{Purpose and corrected scope}

Magnetic helicity is often written as
\[
H=\int_\Sigma \bm A\cdot\bm B\,d^3x
  =\int_\Sigma a\wedge da.
\]
The compact formula hides three separate questions.  First, does a global
potential \(a\) exist?  Second, which gauge transformations and boundary
conditions leave the integral unchanged?  Third, if a projector defines a
moving coherent line, is the relevant field the ambient Abelian field or the
curvature of the induced line connection?

The needed ingredients are standard but belong to different literatures:
Chern--Simons transgression \cite{chernsimons1974}, magnetic and relative
helicity \cite{bergerfield1984,mactaggartvalli2023}, and Berry holonomy
\cite{berry1984,simon1983}.  Keeping their domains distinct is what makes the
combined statement reliable.

The purpose of this paper is to answer those questions in a typed order.  The
answer to the third is especially important.  The projector contribution is
not an extra force added after a Chern--Simons balance has been written.  It is a
part of the induced curvature itself.  Consequently, the exact balance law for
the total induced field has no separate ``projector remainder.''  Such a term
appears only when one compares the total induced field with a surrogate that
has deliberately omitted the Berry contribution.

This distinction also fixes the relation to MTT.  Projection and spectral gaps
can control a supplied coherent line.  They do not, by themselves, prove that
this line is the physical photon sector or that its curvature obeys Maxwell,
magnetohydrodynamic, or other chosen dynamics.  Those identifications belong to
a source contract stated in \cref{sec:mtt}.

\section{The typed geometric contract}

Let \(I=[t_0,t_1]\), let \(\Sigma\) be an oriented Riemannian
three-manifold, and set \(\cU=I\times\Sigma\).  Write
\(\iota_t:\Sigma\hookrightarrow\cU\) for the slice inclusion.  The basic
objects are listed in \cref{tab:contract}.

\begin{table}[ht]
\centering
\small
\begin{tabular}{>{\raggedright\arraybackslash}p{0.18\textwidth}
                >{\raggedright\arraybackslash}p{0.30\textwidth}
                >{\raggedright\arraybackslash}p{0.40\textwidth}}
\toprule
Object & Type & Role \\
\midrule
\(\cE\to\cU\) & Hermitian vector bundle & Ambient carrier \\
\(D\) & Unitary connection on \(\cE\) & Ambient parallel transport \\
\(P\) & \(C^2\) rank-one orthogonal projection & Defines
\(\cL=\operatorname{Ran}P\) \\
\(\nabla^\cL=P D\) & Unitary connection on \(\cL\) & Total induced connection \\
\(f=-iF_{\nabla^\cL}\) & Real two-form on \(\cU\) & Total induced field strength \\
\(b_t=\iota_t^*f\) & Spatial two-form & Magnetic part \\
\(e_t=-\iota_t^*(\iota_{\partial_t}f)\) & Spatial one-form & Electric part \\
Boundary protocol & Closed slice, decay, or relative data & Makes helicity typed \\
\bottomrule
\end{tabular}
\caption{The minimal geometric data.  No physical interpretation is inferred
from these types alone.}
\label{tab:contract}
\end{table}

The sign in the definition of \(e_t\) is chosen so that, in a global gauge
\(\mathcal A=\phi\,dt+a_t\),
\[
f=d\mathcal A
  =dt\wedge(\dot a_t-d_\Sigma\phi)+d_\Sigma a_t,
\qquad
e_t=d_\Sigma\phi-\dot a_t.
\]
Other sign conventions are equivalent after changing all balance formulas
consistently.

\section{The induced line connection}

\begin{definition}[Projected connection]\label{def:projected-connection}
For a section \(s\) of \(\cL=\operatorname{Ran}P\), define
\[
\nabla^\cL s:=P(Ds).
\]
\end{definition}

Because \(P\) is orthogonal and \(D\) is unitary,
\(\nabla^\cL\) is a unitary connection.  It depends on both the ambient
connection and the way the range of \(P\) turns inside \(\cE\).

\begin{theorem}[Curvature of the projected connection]
\label{thm:projected-curvature}
Let \(F_D=D^2\) and let \(DP=[D,P]\).  The curvature of
\(\nabla^\cL\) is
\[
F_{\nabla^\cL}
 =P F_D P+P(DP)\wedge(DP)P
\]
as an endomorphism-valued two-form restricted to \(\cL\).
\end{theorem}

\begin{proof}
For \(s=Ps\),
\[
(P D)^2s=P D(PDs)
       =P(DP)\wedge Ds+P D^2s.
\]
Differentiate \(P^2=P\) to obtain \(P(DP)P=0\).  Splitting
\(Ds=P Ds+(1-P)Ds\) and using
\((1-P)Ds=(DP)s\) gives the stated Grassmann-curvature term.
\end{proof}

The most transparent specialization is an ambient scalar \(U(1)\) connection.
Suppose that, in a local frame,
\[
D=d+iA\,\Id_{\cE},
\]
where \(A\) is a real one-form.  Then \(DP=dP\), \(F_D=i\,dA\), and the
real curvature of the induced line is
\begin{equation}
f=dA+q_P,
\qquad
q_P:=-i\,\Tr\!\bigl(P\,dP\wedge dP\bigr).
\label{eq:curvature-split}
\end{equation}
The second term is the Berry or Grassmann curvature.  In a local unit section
\(u\) with \(P=|u\rangle\langle u|\), the induced real connection form is
\[
\mathcal A=A-i\langle u,du\rangle,
\qquad
f=d\mathcal A.
\]
Under \(u\mapsto e^{i\chi}u\), this local form changes by
\(\mathcal A\mapsto\mathcal A+d\chi\), while \(f\) is unchanged.

\section{When helicity is a real-valued functional}
\label{sec:helicity-domain}

The local formula for \(\mathcal A\) does not guarantee a global real one-form
on a slice.  A nontrivial line bundle has no global unit frame, and even a
trivial bundle requires a boundary and gauge convention before
\(\int a\wedge da\) is an invariant real number.

\begin{assumption}[Absolute-helicity domain]\label{ass:absolute-domain}
For each time under consideration, \(\cL|_{\Sigma_t}\) is supplied with a
global unitary trivialization.  In that trivialization,
\(a_t=\iota_t^*\mathcal A\) and \(b_t=d_\Sigma a_t\).  In addition, one of the
following holds:
\begin{enumerate}[label=(\alph*)]
\item \(\Sigma\) is closed and allowed gauge transformations lift to
single-valued real functions \(\chi\);
\item \(\Sigma=\mathbb R^3\) and all fields and gauge functions decay enough
to remove the boundary term; or
\item a boundary gauge is fixed so that the flux terms displayed below are
part of the data.
\end{enumerate}
\end{assumption}

\par\medskip
\begin{definition}[Induced helicity]\label{def:helicity}
On the domain of Assumption~\ref{ass:absolute-domain}, define
\[
H_\cL(t):=\int_\Sigma a_t\wedge b_t.
\]
\end{definition}

\begin{proposition}[Gauge change]\label{prop:gauge-change}
Under \(a_t\mapsto a_t+d_\Sigma\chi_t\),
\[
H_\cL\mapsto H_\cL+\int_{\partial\Sigma}\chi_t b_t.
\]
Thus \(H_\cL\) is invariant for closed slices and for the decay or fixed
boundary protocols in Assumption~\ref{ass:absolute-domain}.  On a general open or
multiply connected domain, one must instead supply a reference field and use
relative helicity, or retain the boundary term explicitly.
\end{proposition}

\begin{proof}
Since \(d_\Sigma b_t=0\),
\[
\int_\Sigma d_\Sigma\chi_t\wedge b_t
 =\int_\Sigma d_\Sigma(\chi_t b_t)
 =\int_{\partial\Sigma}\chi_t b_t.
\]
\end{proof}

For a nontrivial line bundle, Abelian Chern--Simons data can still be formulated
relative to a reference connection or as a differential character, generally
with a value defined modulo a period lattice.  That construction is not the
same object as the real-valued absolute helicity in
Definition~\ref{def:helicity}.  This
paper uses the latter only on its declared domain.

\section{Exact balance and boundary flux}
\label{sec:balance}

\begin{theorem}[Exact induced-helicity balance]\label{thm:balance}
Assume Assumption~\ref{ass:absolute-domain} and write the full induced connection as
\(\mathcal A=\phi\,dt+a_t\).  Then
\begin{equation}
\frac{d}{dt}H_\cL(t)
=-2\int_\Sigma e_t\wedge b_t
 \int_{\partial\Sigma}\bigl(\phi b_t+a_t\wedge e_t\bigr).
\label{eq:balance}
\end{equation}
In particular, on a closed slice or under a protocol that removes the displayed
boundary flux,
\[
\frac{d}{dt}H_\cL(t)=-2\int_\Sigma e_t\wedge b_t.
\]
The formula contains no additional projector remainder.
\end{theorem}

\begin{proof}
Differentiate \(H_\cL=\int_\Sigma a_t\wedge d_\Sigma a_t\).  Stokes'
theorem gives
\[
\frac{dH_\cL}{dt}
=2\int_\Sigma \dot a_t\wedge b_t
-\int_{\partial\Sigma}a_t\wedge\dot a_t.
\]
Substitute \(\dot a_t=d_\Sigma\phi-e_t\), use
\(d_\Sigma b_t=0\), and integrate the exact boundary two-form
\(d_\Sigma(\phi a_t)\) over \(\partial\Sigma\).  This yields
\cref{eq:balance}.
\end{proof}

\par\medskip
\begin{corollary}[Conservation criterion]\label{cor:conservation}
Under vanishing boundary flux, \(H_\cL\) is conserved whenever
\(\int_\Sigma e_t\wedge b_t=0\).  Pointwise ideal evolution
\(e_t\wedge b_t=0\) is sufficient but not necessary.
\end{corollary}

Small projector derivatives alone do not imply helicity conservation.
They control only the Berry part of the field.  The ambient contribution to
\(e_t\wedge b_t\) may remain nonzero.

\section{What projector control actually bounds}
\label{sec:comparison}

Let
\[
f_0=dA,\qquad f=f_0+q_P
\]
be the ambient and induced real curvatures from
\cref{eq:curvature-split}.  Decompose them on a slice as
\[
b=b_0+\Delta b,\qquad e=e_0+\Delta e,
\]
where
\[
\Delta b=\iota_t^*q_P,
\qquad
\Delta e=-\iota_t^*(\iota_{\partial_t}q_P).
\]

\begin{proposition}[Berry-curvature component bounds]
\label{prop:berry-bounds}
For tangent vectors \(X,Y\),
\[
|q_P(X,Y)|
\le 2\,\|d_XP\|_{\mathrm{op}}\,\|d_YP\|_{\mathrm{op}}.
\]
Consequently, up to fixed norm-equivalence constants determined by the metric
on \(\Sigma\),
\[
|\Delta b|\le 2|d_\Sigma P|^2,
\qquad
|\Delta e|\le 2|\partial_tP|\,|d_\Sigma P|.
\]
\end{proposition}

\begin{proof}
Expand the wedge product:
\[
q_P(X,Y)
=-i\,\Tr\!\left(P(d_XP\,d_YP-d_YP\,d_XP)\right).
\]
Because \(P\) has rank one, the trace of \(PT\) is bounded by
\(\|T\|_{\mathrm{op}}\).  Apply submultiplicativity to the two terms.
\end{proof}

The next result places the old ``remainder'' in its correct role.  It measures
the error in using the ambient field instead of the induced field; it is not an
extra term in \cref{thm:balance}.

\par\medskip
\begin{theorem}[Ambient-surrogate rate error]\label{thm:surrogate-error}
Assume the boundary flux vanishes and all displayed fields are in \(L^2\).
Let
\[
\dot H_{\rm ind}:=-2\int_\Sigma e\wedge b,
\qquad
\dot H_{\rm amb}:=-2\int_\Sigma e_0\wedge b_0.
\]
Then
\begin{align}
|\dot H_{\rm ind}-\dot H_{\rm amb}|
\le 2\bigl(&\|\Delta e\|_{L^2}\|b_0\|_{L^2}
 \|e_0\|_{L^2}\|\Delta b\|_{L^2}\notag\\
&+\|\Delta e\|_{L^2}\|\Delta b\|_{L^2}\bigr).
\label{eq:surrogate-bound}
\end{align}
\end{theorem}

\begin{proof}
Expand
\((e_0+\Delta e)\wedge(b_0+\Delta b)-e_0\wedge b_0\)
and apply Cauchy--Schwarz to each of the three remaining pairings.
\end{proof}

\par\medskip
Combining Proposition~\ref{prop:berry-bounds} and
Theorem~\ref{thm:surrogate-error} gives a quantitative
adiabatic comparison whenever \(\partial_tP\) and \(d_\Sigma P\) are
controlled.  The estimate does not erase the Berry term and does not turn a
nonconserved ambient field into a conserved induced one.

\section{Hopf specialization and normalization}
\label{sec:hopf}

Let \(\Omega_{S^2}\) be the standard oriented area form on the unit
two-sphere, normalized by
\[
\int_{S^2}\Omega_{S^2}=4\pi.
\]

\begin{theorem}[Whitehead integral formula in the chosen normalization]
\label{thm:hopf}
Let \(n:S^3\to S^2\) be smooth.  Since
\(H^2_{\mathrm{dR}}(S^3)=0\), choose a real one-form \(a\) satisfying
\[
da=n^*\Omega_{S^2}.
\]
Then
\[
\operatorname{Hopf}(n)
=\frac{1}{16\pi^2}\int_{S^3}a\wedge da
\in\mathbb Z.
\]
The value is independent of the choice of \(a\).
\end{theorem}

\begin{proof}
This is Whitehead's integral formula with the area form of total area
\(4\pi\).  If \(a'\) is another primitive, then \(a'-a\) is closed and hence
exact on \(S^3\).  The resulting change of the integral is the integral of an
exact three-form and vanishes.
\end{proof}

\par\medskip
If instead one uses \(\omega_{S^2}=\Omega_{S^2}/(4\pi)\), whose integral is
one, and \(d\alpha=n^*\omega_{S^2}\), the same statement reads
\(\operatorname{Hopf}(n)=\int_{S^3}\alpha\wedge d\alpha\).  Mixing these two
normalizations is a common source of an erroneous factor of \(16\pi^2\);
the integral construction is due to Whitehead \cite{whitehead1947}.

For an induced MTT line connection, integer Hopf quantization follows only if
its spatial curvature is actually of the form \(n^*\Omega_{S^2}\) for a
specified map \(n\).  Rank one, coherence, or a spectral gap does not imply
that pullback condition.

\section{Riesz projectors and correctly scaled derivative bounds}
\label{sec:riesz}

Projector control can be obtained from a gapped operator family, but the
operator domain and the gap scaling must be explicit.

\begin{theorem}[Bounded-family Riesz estimate]\label{thm:riesz}
Let \(y\mapsto L(y)\) be a \(C^1\) family of bounded self-adjoint operators on
a fixed Hilbert space.  Suppose
\[
\sigma(L(y))\subset\{0\}\cup[\lambda_*,\infty),
\qquad \lambda_*>0,
\]
and the multiplicity of the zero eigenspace is constant.  Let \(P(y)\) be the
orthogonal projection onto \(\ker L(y)\).  For the positively oriented circle
\(\Gamma=\{z:|z|=\lambda_*/2\}\),
\[
P(y)=\frac{1}{2\pi i}\oint_\Gamma (z-L(y))^{-1}\,dz
\]
and, for every parameter direction \(v\),
\[
\partial_vP
=\frac{1}{2\pi i}\oint_\Gamma
(z-L)^{-1}(\partial_vL)(z-L)^{-1}\,dz.
\]
Moreover,
\begin{equation}
\|\partial_vP\|_{\mathrm{op}}
\le \frac{2}{\lambda_*}\,
\|\partial_vL\|_{\mathrm{op}}.
\label{eq:riesz-bound}
\end{equation}
\end{theorem}

\begin{proof}
Differentiate the resolvent identity under the contour integral.  On
\(\Gamma\), self-adjointness and the spectral assumption give
\(\|(z-L)^{-1}\|\le2/\lambda_*\).  Since
\(\operatorname{len}(\Gamma)=\pi\lambda_*\), the contour estimate is
\[
\frac{\pi\lambda_*}{2\pi}
\left(\frac{2}{\lambda_*}\right)^2
\|\partial_vL\|
=\frac{2}{\lambda_*}\|\partial_vL\|.
\]
\end{proof}

For unbounded Laplace-type families, norm differentiability in
\(\mathcal B(\calH)\) is generally the wrong hypothesis.  One instead fixes a
common domain or uses graph-norm and relative-resolvent bounds, as in Kato
perturbation theory.  The same contour identity remains available after those
domain hypotheses are supplied.  This paper does not silently replace that
unbounded problem by the bounded theorem above.  The relevant perturbation
framework is developed systematically by Kato \cite{kato1995}.

\section{Conditional MTT encoding}
\label{sec:mtt}

The preceding results are ordinary differential geometry and functional
analysis.  MTT enters only through a proposed source for the typed data.  The
current scope is aligned with the revised Foundation, fixed-point, and
topological-encoding papers
\cite{nero_foundations_v9,nero_fixedpoints_i_v7,nero_topological_v2}.

\begin{definition}[Electromagnetic source contract]\label{def:source-contract}
An MTT electromagnetic source packet on \(\cU\) must provide:
\begin{enumerate}[label=\textnormal{S\arabic*.},leftmargin=3.5em]
\item a Hermitian carrier \((\cE,D)\);
\item a \(C^2\), constant-rank, rank-one projector \(P\);
\item a theorem identifying \(\operatorname{Ran}P\), rather than merely an
isomorphic line, with the physical electromagnetic sector;
\item an intertwiner identifying \(-iF_{PD}\) with the physical
electromagnetic field strength and preserving its dynamics;
\item a trivialization, relative connection, or differential-character
protocol that types the claimed helicity observable;
\item when Hopf quantization is claimed, a selected map
\(n:S^3\to S^2\) and the equality
\(\iota_t^*(-iF_{PD})=n^*\Omega_{S^2}\).
\end{enumerate}
\end{definition}

\begin{theorem}[Conditional pullback of helicity statements]
\label{thm:conditional-pullback}
If an MTT source packet satisfies S1--S5, then
Theorem~\ref{thm:projected-curvature},
Proposition~\ref{prop:gauge-change}, Theorem~\ref{thm:balance},
Proposition~\ref{prop:berry-bounds}, and
Theorem~\ref{thm:surrogate-error} apply to the emitted physical
electromagnetic sector.  If S6 also holds, \cref{thm:hopf} supplies the
integer Hopf invariant.  If \(P\) is emitted as the kernel projector of a
bounded family satisfying \cref{thm:riesz}, then
\cref{eq:riesz-bound} controls its parameter derivatives.
\end{theorem}

\begin{proof}
Each source row identifies an emitted object with the corresponding typed
hypothesis of the cited result.  Composition with the stated intertwiners
therefore transports those identities without changing their domains.
\end{proof}

\begin{table}[ht]
\centering
\small
\begin{tabular}{>{\raggedright\arraybackslash}p{0.12\textwidth}
                >{\raggedright\arraybackslash}p{0.31\textwidth}
                >{\raggedright\arraybackslash}p{0.43\textwidth}}
\toprule
Rows & Status here & Meaning \\
\midrule
S1--S2 & Conditional representation data & Standard projected-connection
mathematics applies once supplied \\
S3--S4 & Open source identification & No current theorem in this paper selects
the physical photon line and its dynamics \\
S5 & Case-dependent & Closed, decaying, relative, and differential-character
domains are distinct \\
S6 & Specialization only & Hopf quantization requires the explicit pullback
condition \\
\bottomrule
\end{tabular}
\caption{Claim status of the MTT bridge.}
\label{tab:mtt-status}
\end{table}

The conditional theorem is useful because it states exactly what remains to be
constructed.  It is not evidence that S3--S4 already hold.

\section{What the paper establishes}

\vspace{0.4\baselineskip}
The durable mathematical conclusions are:
\begin{enumerate}
\item a rank-one projector and ambient unitary connection define an induced
line connection with the exact curvature in
\cref{thm:projected-curvature};
\item absolute helicity is a typed observable only after trivialization,
gauge, and boundary data are fixed;
\item the exact balance law is \cref{eq:balance};
\item omission of the Berry curvature has the controlled comparison error in
\cref{eq:surrogate-bound};
\item Hopf quantization is valid under the explicit pullback and normalization
in \cref{thm:hopf}; and
\item a spectral gap controls a bounded-family Riesz projector with scaling
\(1/\lambda_*\), not an unexplained \(1/\lambda_*^2\) factor.
\end{enumerate}

The paper does not derive Maxwell's equations, ideal magnetohydrodynamics, a
physical photon sector, a universal helicity conservation law, or a universal
topological classification from MTT.  Those claims require the source and
dynamical rows in Definition~\ref{def:source-contract}.

\section{Version 3 changes and reasons}

\begin{center}
\small
\begin{tabular}{>{\raggedright\arraybackslash}p{0.25\textwidth}
                >{\raggedright\arraybackslash}p{0.31\textwidth}
                >{\raggedright\arraybackslash}p{0.32\textwidth}}
\toprule
Change & Earlier issue & Reason \\
\midrule
Retitled and reclassified as a conditional encoding &
The previous title implied that MTT had selected the physical coherent
electromagnetic line &
The selected physical source remains open \\
\addlinespace
Replaced the projector-remainder balance &
Projector variation was counted once in the induced field and again as a
remainder &
The exact balance must use the full connection; comparison error is a separate
question \\
\addlinespace
Separated spacetime and slice forms &
The electric field was defined by contracting a spatial two-form &
The contraction must be taken before pullback to the slice \\
\addlinespace
Added gauge, topology, and boundary domains &
A local connection form was treated as a global helicity potential &
Nontrivial bundles and open domains require relative or differential data \\
\addlinespace
Fixed the Hopf convention &
The normalized area form and the factor \(1/(16\pi^2)\) were mixed &
The coefficient depends on whether the sphere area is \(4\pi\) or \(1\) \\
\addlinespace
Corrected the Riesz estimate &
The contour constant obscured the gap dimension and unbounded-domain issue &
The bounded theorem gives \(2\|\partial L\|/\lambda_*\); Laplace-type families
need graph-domain hypotheses \\
\bottomrule
\end{tabular}
\end{center}

\section{Discussion}

The main conceptual lesson is modest but useful.  A varying projector can
produce a genuine geometric contribution to an Abelian field strength.  That
contribution is neither optional nor mysterious: it is the curvature of the
projected connection.  Once it is included, the usual Chern--Simons calculus
works without modification.  The new quantitative question is how far the
induced field lies from an ambient-field approximation, and
\cref{thm:surrogate-error} answers that question on its stated domain.

For MTT, this yields a clean research target.  A future source theorem should
not merely point to a rank-one coherent subspace.  It must identify the same
line and connection with the physical electromagnetic sector, transport the
dynamics, and specify the global observable domain.  If that target is met,
the geometric results in this paper apply immediately.  Until then, they are a
rigorous compatibility and completion contract.

\bibliographystyle{plain}
\bibliography{main}

\end{document}
